\documentclass[11pt, a4paper]{article}

\usepackage[top = 0.6 in, bottom = 0.6 in, left = 0.8 in, right = 0.8 in]{geometry}
\usepackage{amsmath, amssymb, amsfonts}

\allowdisplaybreaks[4]

\usepackage{enumerate}
\usepackage{array}
\usepackage{multirow}
\usepackage{dingbat}
\usepackage{fontawesome5}
\usepackage{tasks}
\usepackage{bbding}
\usepackage{undertilde}
\usepackage{twemojis}
% how to use bull's eye ----- \scalebox{2.0}{\twemoji{bullseye}}
\usepackage{fontspec}
\usepackage{fancyhdr}
\usepackage{lastpage}

\pagestyle{fancy}
\fancyhf{}
\fancyfoot[C]{Page \thepage\ of \pageref{LastPage}}
\renewcommand{\headrulewidth}{0pt}

\title{MSMS 402 : Assignment 01}
\author{{\fontsize{25}{20} \benyorithafont Ananda Biswas} \\[1em] Exam Roll No. : 24419STC053}
\date{\today}

\newfontface\benyorithafont{Benyoritha-G334D.otf}

\newfontface\myfont{Myfont1-Regular.ttf}[LetterSpace=0.05em]

\newfontface\cbfont{CaveatBrush-Regular.ttf}
% how to use --- \myfont --write text here--

\newfontface\gloriahallelujahfont{GLORIAHALLELUJAH-REGULAR.TTF}

\begin{document}

\maketitle

\textbf{Theorem.} The acceptance--rejection algorithm generates a random variable $X$ such that
\begin{align*}
P\{X = j\} = p_j, \quad j = 0, \ldots
\end{align*}
In addition, the number of iterations of the algorithm needed to obtain $X$ is a geometric random variable with mean $c$. \\[1em]

\textbf{Proof.} As a proposal distribution, we use random variable $Y$ such that
\begin{align*}
P\{Y = j\} = q_j, \quad j = 0, \ldots.
\end{align*}

To begin, let us determine the probability that a single iteration produces the accepted value $j$. First note that
\begin{align*}
P\{Y = j, \text{ it is accepted}\}
&= P\{Y = j\} P\{\text{accept} \mid Y = j\} \\[0.5em]
&= q_j \dfrac{p_j}{c q_j} \\[0.5em]
&= \dfrac{p_j}{c}.
\end{align*}

Summing over $j$ yields the probability that a generated random variable is accepted:
\begin{align*}
P\{\text{accepted}\}
&= \sum_j \dfrac{p_j}{c} \\[0.5em]
&= \dfrac{1}{c}.
\end{align*}

As each iteration independently results in an accepted value with probability $\dfrac{1}{c}$, we see that the number of iterations needed is geometric with mean $c$. Also,
\begin{align*}
P\{X = j\}
&= \sum_n P\{\text{$j$ accepted on iteration $n$}\} \\[0.5em]
&= \sum_n \left(1 - \dfrac{1}{c}\right)^{n-1} \dfrac{p_j}{c} \\[0.5em]
&= p_j.
\end{align*}

\flushright{\Box}



\end{document}